\documentclass{article}
\usepackage{decomp}
\begin{document}
\decompnote{Collection markers: a cited work that is in the local collection is tagged \localcopy{<key>} after its reference entry; it is decompiled in decompiled/<key>/. A reference to another version of a local work (e.g. its preprint) is tagged \localother{<key>}. A reference without that tag is not in the local collection and needs an external pull. Tags added by \texttt{\detokenize{tools/mark_local.py}}.}

\sourcefile{10.1515\_mim-2026-0015.pdf}
\decompnote{Publisher's version of record (De Gruyter Brill, Methods in Microscopy, typeset PDF, 10 pages, published online ahead of print, ``aop''). Running headers (``Y. Hua et al.: Development of sample preparation for volume electron microscopy'', ``DE GRUYTER'') and page numbers, the De Gruyter logo and the open-access icon are left out. The journal line at the top of p.\,1 and the licence footer of p.\,1 are kept. The journal's typesetting puts the figures at the top of pages; here each figure follows the paragraph that is interrupted by it.}

\pdfpage{1}
\noindent Methods in Microscopy 2026; aop\par

\noindent\textbf{Review Article}\par

\bigskip
\noindent Yunfeng Hua*, Linlin Li* and Hua Han*\par

\medskip
{\raggedright\noindent\LARGE\textbf{Preparing large biological specimens for high-throughput volume electron microscopy}\par}

\bigskip
\decompnote{The corresponding-author footnote below is printed at the foot of the left column of p.\,1; the licence statement at the foot of p.\,1 (between ``Thoss Media'' and ``This work'' a CC BY licence badge is printed).}
{\small\noindent \textbf{*Corresponding authors:} \textbf{Yunfeng Hua}, Shanghai Institute of Precision Medicine, Ninth People’s Hospital, Shanghai Jiao Tong University School of Medicine, 200125, Shanghai, China, E-mail: yunfeng.hua@shsmu.edu.cn; \textbf{Linlin Li}, State Key Laboratory of Brain Cognition and Brain-inspired Intelligence Technology, Institute of Automation, Chinese Academy of Sciences, 100190, Beijing, China, E-mail: lilinlin@ia.ac.cn; and \textbf{Hua Han}, State Key Laboratory of Brain Cognition and Brain-inspired Intelligence Technology, Institute of Automation, Chinese Academy of Sciences, 100190, Beijing, China; and School of Future Technology, University of Chinese Academy of Sciences, 101408, Beijing, China, E-mail: hua.han@ia.ac.cn\par}

{\small\noindent Open Access. © 2026 the author(s), published by De Gruyter on behalf of Thoss Media This work is licensed under the Creative Commons Attribution 4.0 International License.\par}

\medskip
{\small\noindent \url{https://doi.org/10.1515/mim-2026-0015}\\
Received July 30, 2026; accepted August 31, 2026;\\
published online September 24, 2026\par}

\medskip
\noindent \textbf{Abstract:} To enable high-throughput volume electron microscopy (vEM) imaging, osmium-based staining methods are routinely employed to enhance membrane contrast and electrical conductivity in biological specimens. Nevertheless, reliable preparation of densely and uniformly stained tissue blocks remains a major challenge for large-scale vEM-based ultrastructure analysis, especially for connectomic studies. Recently, multiple strategies have been proposed to achieve high-quality \emph{en bloc} staining of millimeter-sized brain samples without compromising membrane contrast, paving the way for system-level neural circuit mapping using vEM. Herein, we review the development of large-volume tissue staining and elaborate the underlying chemical logic based on cumulative experimental evidence. We further discuss emerging challenges associated with adapting current staining strategies for cryo-preserved tissue samples and in-resin applications of correlative light and electron microscopy.

\medskip
\noindent \textbf{Keywords:} volume electron microscopy; large-scale connectomics; osmium staining; \emph{en bloc} tissue processing; high-throughput imaging

\bigskip
\noindent Over the past two decades, conventional serial ultra-thin section electron microscopy has evolved into a group of emerging techniques that enable automated acquisition of large-scale volumetric image datasets from samples thicker than one micron, collectively termed volume electron microscopy (vEM) [1]. Based on imaging principles, vEM can be classified into two major categories: serial section-based and serial block-face-based methods [2]–[5]. Serial section-based vEM implements automated collection of ultrathin sections either on opaque tape for scanning electron microscopy (SEM) [6], or on grid tape for transmission electron microscopy (TEM) [7] and electron tomography (ET) [8]. Since tape-related artifacts complicate downstream data analysis, direct silicon wafer-based section collection techniques have recently been developed with the assistance of magnetic field-based systems (MagC [9], GAUSS-EM [10]) or an ionizer-equipped device (Drift-EM [11]). These methods are non-destructive and therefore allow post-staining of mounted sections, but additional sample processing inevitably causes section loss, folding, contamination, especially for collecting large datasets. In contrast, serial block-face vEM produce intrinsically aligned image stacks by iteratively imaging newly exposed sample surfaces through in-chamber ultramicrotomy cutting [12] or focused ion beam milling [13]. While this workflow circumvents the major drawbacks of traditional ultra-thin sectioning, it mandates the use of \emph{en bloc}-stained tissue blocks due to an integrated imaging and surface removal process. This scenario also applies to large-scale ion beam etching of thick sections [14], [15], which is regarded as a parallelized block-face imaging modality. Furthermore, sufficient heavy-metal impregnation is indispensable for mitigating charging artifact under SEM, especially for imaging large biological samples or conducting fast imaging under high electron dose. As supplementary anti-charging strategies, conductive resin embedding [16], carbon nanotube-coated ATUM tape [17], and in-chamber focal charge compensation [18] are often required to image specific tissue samples prone to charging.

\section*{1 Development of staining protocol for vEM}

\noindent Driven by the ambitious goal of imaging intact neural circuits at nanometer resolution in large tissue volumes [19], there remains an ongoing debate regarding the \pdfpage{2}ultimate imaging method of choice. As the best (if not the only) option for serial block-face vEM users, the combination of ferrocyanide-reduced osmium [20] with thiocarbohydrazide (TCH)-mediated osmium amplification [21] – termed rOTO staining [22] – yields tissue specimens with superior membrane contrast and adequate electrical conductivity [23]–[25]. When adapting the original rOTO protocol for mouse cortical tissues routinely collected via 1 mm-diameter biopsy needles (Figure 1A) [26] as well as whole mouse brain [27], Hua and Mikula found independently that homogenous, high-quality staining was only achievable within a superficial layer of approximately 200 μm from the tissue surface. Note that these stained samples exhibited low-contrast central regions encapsulated by a dense layer of black precipitates (Figure 1B). Such staining inhomogeneity could not be simply attributed to slow reagent penetration and persisted even after prolonged incubation, representing a substantial bottleneck for \emph{en bloc} staining of large-volume tissues. This notion was later supported by real-time monitoring of the stain-penetration process using X-ray microscopy [28]. In addition, a 10–15 \% sample shrinkage in aldehyde-fixed brain tissues was observed after treatment with osmium and ferrocyanide mixtures, offering one plausible explanation for the observed diffusion-barrier formation [28].

\begin{center}
\includegraphics[width=0.8\linewidth,height=0.55\textheight,keepaspectratio]{figures/fig1.png}
\end{center}
\noindent \textbf{Figure 1:} Increasing stain penetration depth for 1-mm-sized brain samples (modified from ref. [26]). (A) workflow for screening 1-mm-sized brain biopsy samples. Samples were \emph{en bloc} stained and embedded in resin. A block face was prepared by cutting the sample in half, which was then imaged for screening in overview (1), peripheral (2), intermediate (3), and core regions (4) under a scanning electron microscope. (B, C) Representative images from the corresponding regions in A (regions 1–4, left to right columns) for mouse cortex samples stained by rOTO (B) and ORTO (C). ORTO eliminates the precipitate ring that frequently forms in rOTO-stained samples.
\alttext{Figure 1 is printed at the top of p.\,2, across about two thirds of the page width. Three rows. (A) Schematic, black and grey line drawing, left to right with arrows between the steps: an open cylinder labelled ``Biopsy'' with dimension arrows ``1 mm'' (diameter, top) and ``>1 mm'' (height, left side, written vertically); an arrow to a dark grey cylinder labelled ``En bloc staining''; an arrow to the cylinder cut into two short pieces, one showing a white cut face, labelled ``Block-face imaging''; an arrow to a grey disc (the block face) with a large black square frame labelled ``1'' covering its upper left part and three small squares labelled ``2'', ``3'' and ``4'' running from near the disc's upper edge inward towards its centre. (B) Row of four greyscale SEM images labelled ``B'' at upper left, each with a white boxed number 1--4 in the middle; the row is labelled ``rOTO'' below. B1 (overview): the corner of the round sample, a pale band of embedding resin and a dark outer rim at upper left, then grey tissue with many small pale dots; a dark curved band runs through the tissue roughly parallel to the edge, from the bottom edge (marked by a black arrowhead below the image, labelled ``Precipitate ring'') up towards the upper right; inside the band the tissue is a somewhat darker grey. White scale bar ``200 μm''. B2 (periphery): neuropil with clearly outlined dark membranes, round and oval profiles and a dark-ringed profile at upper left; white scale bar ``2 μm''. B3 (intermediate) and B4 (core): low-contrast, mottled grey tissue in which membranes are faint and blurred, a few dark-ringed profiles and pale patches; B3 has a dark diagonal band along its left side. (C) Same layout for ORTO, labelled ``C'' at upper left and ``ORTO (Hua protocol)'' below. C1: the sample corner with dark outer rim and evenly grey tissue with small pale dots, without an internal dark band. C2--C4: neuropil with sharp dark membranes at similar contrast in all three; thick black rings and bands (myelin-like profiles) in each, a bundle of thick black streaks running diagonally through C2, a large dark ring at lower left of C3 and at right of C4. No scale bars in row C.}

As the practical utility of serial block-face vEM is constrained by the maximum sample size processable using rOTO staining, early methodological advances in this field were primarily pioneered by Winfried Denk and his followers (Table 1). In 2015, three independent research teams proposed distinct strategies to resolve the bottleneck of large-volume sample preparation. Collectively, these studies demonstrated that the mixture of OsO\textsubscript{4} and ferrocyanide (“reduced osmium”) is responsible for the undesirable precipitation barrier that forms within brain tissue samples [26], [27]. As a straightforward solution, Hua et al. developed a modified procedure featuring sequential OsO\textsubscript{4} penetration and ferrocyanide reduction to replace the conventional reduced-osmium step in the classic rOTO protocol [26]. Note that an intermediate wash step between OsO\textsubscript{4} and ferrocyanide incubation is considered unnecessary, as reduced-osmium species will be generated to some extent. This simple modification, termed the ORTO protocol (also informally referred to as Hua protocol), enables uniform staining of biopsy samples up to 1 mm in diameter (Figure 1C). In parallel, Mikula and Denk established the BROPA (brain-wide reduced-osmium staining with pyrogallol-mediated amplification, or rOPO) protocol [27]. This procedure incorporates 10 \% formamide to stabilize reduced-osmium species and prolong their efficacy during extended incubation. Besides, \pdfpage{3}pyrogallol was used for the first time to replace TCH as a better-suited reagent for osmium amplification, supporting a 10-day-long incubation step required for whole mouse brain samples [27]. Notably, a significantly shortened BROPA protocol (fBROPA) also enables homogeneous staining of millimeter-sized fish brain samples within 4 days [35]. Moreover, Palloto et al. showed optimized preservation of the extracellular space (ECS) by elevating osmolarity during tissue fixation [29], which physically facilitates deeper stain diffusion and thereby supports uniform staining of large-volume samples [27], [29], [36].

\decompnote{Table 1 is printed rotated by 90° on the left half of p.\,3. The two row labels ``rO'' and ``OTO'' are printed on the first line of the two- and three-line blocks they belong to.}
\noindent\textbf{Table 1:} Summary of reduced osmium-based staining methods.

\begin{center}
\small
\resizebox{\linewidth}{!}{%
\begin{tabular}{llllllll}
\toprule
 & \multicolumn{7}{c}{\textbf{Published protocols}}\\
\cmidrule{2-8}
 & \textbf{Deerinck et al. [24]} & \textbf{Mikula et al. [27]} & \textbf{Hua et al. [26]} & \textbf{Pulloto et al. [29]} & \textbf{Polilov et al. [30]\textsuperscript{c}} & \textbf{Song et al. [31]\textsuperscript{c}} & \textbf{Karlupia et al. [32]}\\
\midrule
Fixation & PFA + GA & GA + sucrose & PFA + GA & PFA + GA (high Osm) & GA + OsO\textsubscript{4} & PFA + GA & PFA + GA (high Osm)\\
rO & & & OsO\textsubscript{4} & & OsO\textsubscript{4} & OsO\textsubscript{4} & OsO\textsubscript{4}\\
 & OsO\textsubscript{4} + FeCN & OsO\textsubscript{4} + FeCN + 10 \% formamide & FeCN\textsuperscript{a} & OsO\textsubscript{4} + FeCN & FeCN\textsuperscript{a} & FeCN & FeCN\\
OTO & & OsO\textsubscript{4} & (OsO\textsubscript{4})\textsuperscript{b} & & & OsO\textsubscript{4} & \\
 & TCH & Pyrogallol & TCH & TCH & & Pyrogallol & TCH\\
 & OsO\textsubscript{4} & OsO\textsubscript{4} & OsO\textsubscript{4} & OsO\textsubscript{4} & & OsO\textsubscript{4} & OsO\textsubscript{4}\\
\bottomrule
\end{tabular}}
\end{center}
{\footnotesize\noindent GA, glutaraldehyde; Osm, osmolarity; FeCN, ferrocyanide. \textsuperscript{a}Intermediate wash step is not applied. \textsuperscript{b}Optional. Additional OsO\textsubscript{4} step was used to enhance TCH reactivity [33], [34]. \textsuperscript{c}All steps at 4 °C. Necessary for improved tissue integrity.\par}
\decompnote{The table uses ``Pulloto et al. [29]''; reference [29] is M. Pallotto et al., and the text says ``Palloto et al.''. The last sentence of the footnote, ``Necessary for improved tissue integrity.'', is printed on its own line after footnote c; it is not marked with a letter, so which entry it refers to is not stated.}
\decompnote{Also as \texttt{tables/table1.csv}, transcribed from the page image (the PDF text layer drops the subscript 4 of OsO\textsubscript{4} in the table body).}

Since then, \emph{en bloc} staining using these rOTO variants has become popular within the vEM community [14], [37]–[41], including serial-section vEM laboratories, owing to its operational simplicity (eliminating the need for post-staining of sections) and enhanced imaging throughput [7], [42]–[46]. More importantly, different concepts established in these studies provide a foundational framework for developing a common procedure scalable to diverse tissue dimensions (Figure 2). For instance, sequential OsO\textsubscript{4} penetration and ferrocyanide reduction procedure were prolonged and performed at 4 °C to prepare the whole insect brains [30], [47]. More recently, the modified ORTO protocol has been combined with ECS preservation to enable staining of human and rodent surgical brain samples up to 37 mm\textsuperscript{3} in volume, with one dimension below 2 mm [32]. Alternatively, it has also incorporated the TCH-to-pyrogallol substitution (from BROPA) in a prolonged, low-temperature implementation adapted for whole mouse brains and human brain specimens [31].

In addition to compounds of osmium, other heavy metals are adopted to further increase the overall electron contrast of biological structures. A two-step uranyl acetate incubation procedure has been established and implemented after osmium staining of brain tissues [26]. Prolonged incubation (overnight to 2 days) at 4 °C enables long-distance stain diffusion without inducing cytoplasmic darkening, while a short warming step (2–5 h) at room temperature or 50 °C enhances the reactivity of uranyl acetate with protein and membrane targets [26], [31], [32]. For oversized samples, note that the application of a high-concentration (4 \%) staining solution is recommended [31]; alternatively, repeated replacement with fresh solution during incubation can ensure sufficient stain supply. Regarding another widely used staining buffer, the Walton solution (lead aspartate), early efforts have attempted to attenuate stain deposition by lowering the pH of the staining buffer from 5.5 to 5 [26]. However, lead aspartate solution is found unsuitable for long-term incubation due to its high susceptibility to light and air exposure. Consequently, it is not surprisingly that \pdfpage{4}lead staining has been excluded from up-to-date protocols optimized for large-volume tissue samples [31], [32].

\begin{center}
\includegraphics[width=0.95\linewidth,height=0.55\textheight,keepaspectratio]{figures/fig2.png}
\end{center}
\noindent \textbf{Figure 2:} Protocol development for high-contrast \emph{en bloc} staining. Joint efforts have greatly promoted the development of \emph{en bloc} staining techniques for specimens ranging from small brain tissues to whole brains. To the best of our knowledge, Deerinck et al. first implemented rOTO staining for brain samples, laying a foundational framework for subsequent protocol optimization. Hua et al. further proposed a separate OsO\textsubscript{4} penetration and reduction routine (OR, magenta) to achieve homogenous staining of 1-mm samples. Meanwhile, Mikula et al. established a pyrogallol-mediated osmium amplification (PO) strategy and validated that the preservation of extracellular space (ESC, blue), a method originally developed by Palloto et al., can substantially facilitate stain penetration into brain tissues. Notably, 10 \% formamide is supplemented to the reduced osmium mixture (rO) to enhance solution stability during prolonged incubation. Alternatively, Polilov et al. performed a prolonged OR treatment at low temperature (4 °C, yellow) for insect brains. More recently, Song et al. integrated prolonged OR treatment with PO amplification at 4 °C, enabling uniform staining of the entire mouse brain. In parallel, Karlupia et al. optimized room-temperature ORTO staining for surgical samples with preserved ECS. Collectively, the newly established low-temperature ORPO workflow exhibits great potential for further improving the overall efficiency and quality of \emph{en bloc} brain staining.
\alttext{Figure 2 is printed at the top of p.\,4, across about two thirds of the page width. Flow diagram from left to right. Each protocol is a rectangular box with a grey-scale gradient fill (suggesting the staining profile across a sample) and a coloured name; the citation is written below it. Left: box ``rOTO'' (white text; fill with a black band towards the left fading to white at the right), ``Deerinck et al. 2010''. Three black arrows fan out from it to: (top) a box ``rOPO'' (``rO'' white and underlined, ``PO'' red; fill grey, darker at left, lighter at right), ``Mikula et al. 2015''; (middle) a box ``ORTO'' (``OR'' magenta, ``TO'' white; fill almost uniformly dark grey), ``Hua et al. 2015''; (bottom) a box ``rOTO (ECS)'' (``(ECS)'' light blue; grey fill with a black band near the right end and a white strip at the far right), ``Pulloto et al. 2015''. A red arrow goes from rOPO right to a box ``ORPO (4°C )'' (``OR'' magenta, ``PO'' red, ``(4°C )'' yellow; dark fill), ``Song et al. 2023''. A magenta line rises from the ORTO box, passes behind a smaller box ``OR (4°C )'' (``OR'' magenta, ``(4°C )'' yellow; dark fill), ``Polilov et al. 2021'', and leaves its upper right corner as a magenta arrow into the Song box. A second magenta arrow goes from ORTO down to a box ``ORTO (ECS)'' (``OR'' magenta, ``(ECS)'' light blue; dark uniform fill), ``Karlupia et al. 2023''; a light-blue arrow goes from the rOTO (ECS) box right to the same Karlupia box. At the far right a black brain icon receives a red and a yellow arrow from the Song box and a light-blue and a magenta arrow from the Karlupia box. The figure prints ``Pulloto'' and ``ECS''; the caption prints ``Palloto'' and ``ESC''.}
\decompnote{The caption writes ``ESC'' for the extracellular space (``ECS'' in the text, the figure and the next sentences) and ``Palloto et al.''; the figure and Table 1 print ``Pulloto'', reference [29] is ``M. Pallotto''.}

Although these protocols were originally developed for EM imaging, they have gained increasing recognition as ready-to-use sample preparation approaches for X-ray tomography [48], [49] and photoemission electron microscopy [50]. Moreover, osmium-based staining is lipid-selective and in principle readily applicable to other organ types. However, the incubation durations for key steps must still be determined empirically, due to heterogeneous tissue densities. This process can be both costly and time-consuming, as it requires repeated test runs with \emph{in-situ} monitoring via X-ray microscopy and final validation of sample quality using electron microscopy [28], [31].

\section*{2 Putative chemical mechanisms of reduced osmium staining}

\noindent Despite the widespread application of these reduced-osmium-based protocols for preparing various biological specimens, the detailed chemical reactions underlying membrane contrast generation remains incompletely elucidated. A prevailing model describing reduced-osmium staining was initially proposed by Hua et al. [26] and has recently been revisited through chemical analysis [51]. The model can be simplified (Figure 3) and used to rationalize the enhanced membrane-specific osmium deposition achieved by supplementing ferrocyanide, which in consensus functions as a reducing agent [20], [26], [51]. First, OsO\textsubscript{4} exists in both nonpolar (lipophilic) and \pdfpage{5}ionized (water-soluble) forms. This unique property enables efficient penetration of OsO\textsubscript{4} molecules across numerous membrane-water interfaces within thick tissue samples. Second, OsO\textsubscript{4} undergoes a multi-step reduction process that yields water-insoluble OsO\textsubscript{2}. This black solid preferentially deposits on membrane surface, acting as the core substance responsible for generating membrane contrast under electron microscopy [26], [51]. Third, ferrocyanide may directly reduce ionized OsO\textsubscript{4} to intermediate Os\textsuperscript{VI} species outside lipid bilayers, prompting further reduction to OsO\textsubscript{2}. Effectively, the addition of ferrocyanide to the OsO\textsubscript{4} solution establishes a shortcut for OsO\textsubscript{2} generation, thereby maximizing membrane contrasting. However, replacing ferrocyanide with other reducing agents such as Na\textsubscript{2}SO\textsubscript{3}/NaHSO\textsubscript{3} has failed to yield comparable results, implying that ferrocyanide may exert additional roles beyond osmium reduction. Moreover, these proposed reactions are strongly pH-dependent (Figure 3), which warrants further investigations to identify the optimal working conditions for osmium deposition in a controlled fashion.

\begin{center}
\includegraphics[width=0.6\linewidth,height=0.55\textheight,keepaspectratio]{figures/fig3.png}
\end{center}
\noindent \textbf{Figure 3:} Schematic illustration of osmium chemistry involved in membrane staining. Top: OsO\textsubscript{4} molecule (Os\textsuperscript{VIII}) may rapidly penetrate multiple membrane bilayers (dashed arrows) via partitioning between organic and aqueous phases (1). Gradually, membrane-residing Os\textsuperscript{VIII} is reduced to Os\textsuperscript{VI} by unsaturated lipids within the organic phase, followed by subsequent hydrolysis (thin solid arrows) (2). In addition, Os\textsuperscript{VIII} reduction can also occur in water through a ferrocyanide (Fe\textsuperscript{II})-assisted pathway (thick solid arrow) (3). The final contrast agent, OsO\textsubscript{2} (Os\textsuperscript{IV}), is ultimately generated via further reduction of Os\textsuperscript{VI} species either through its disproportionation (4) or reduction by ethanol during the subsequent dehydration procedure (bent arrow). Bottom: putative chemical reactions involved in the reduced osmium staining.
\decompnote{Figure 3 is printed in the right column of p.\,4, beside the text of section 2; it is placed here because the paragraph that starts on p.\,4 continues on p.\,5.}
\alttext{Schematic with a reaction list below it. Top part: a grey arrow pointing right across the top labelled ``Membrane contrasting''; a grey arrow pointing down along the left side labelled ``Osmium penetration'' (vertical text). A lipid bilayer is drawn horizontally (two rows of light-blue ball heads with double-line tails); dashed lines at its two surfaces extend to the left. Vertical labels on the left: ``Aqueous'' (above the bilayer), ``Organic'' (inside), ``Aqueous'' (below). Above the bilayer: a yellow-outlined open box ``Os\textsuperscript{VIII}'' at left; a green-outlined hexagon ``Fe\textsuperscript{II}'' above the gap to its right; a thick black arrow labelled ``(3)'' from the Os\textsuperscript{VIII} box right to a magenta-outlined open box ``Os\textsuperscript{VI}''; from this box a thin black arrow labelled ``(4)'' runs right and then bends down to a black filled box ``Os\textsuperscript{IV}'' (white text) that sits on the upper row of lipid heads. Inside the bilayer: a yellow filled box ``Os\textsuperscript{VIII}'' at left, a thin black arrow right to a pale pink filled box ``Os\textsuperscript{VI}'' labelled ``(2)'', and a thin black arrow up from this box through the upper lipid heads to the Os\textsuperscript{VI} open box above. A dashed black arrow goes down from the upper Os\textsuperscript{VIII} box into the yellow Os\textsuperscript{VIII} box in the membrane, and another dashed arrow continues down through the lower lipid heads to a yellow-outlined open box ``Os\textsuperscript{VIII}'' labelled ``(1)'' in the lower aqueous phase. Bottom part, four reactions; species are highlighted the same way (yellow fill = membrane Os\textsuperscript{VIII}, yellow outline = aqueous Os\textsuperscript{VIII}, green outline = ferrocyanide, magenta outline = Os\textsuperscript{VI}, black fill = Os\textsuperscript{IV}): (1): [yellow fill] Os\textsuperscript{VIII}O\textsubscript{4} + 2OH\textsuperscript{-} ⇄ [yellow outline] [Os\textsuperscript{VIII}O\textsubscript{4}(OH)\textsubscript{2}]\textsuperscript{2-}. (2): an alkene drawn as R–CH=CH–R + [yellow fill] Os\textsuperscript{VIII}O\textsubscript{4} + 2H\textsubscript{2}O + 2OH\textsuperscript{-} ⟶ a diol drawn as R–CH(OH)–CH(OH)–R (written ``CH-CH'' with HO and OH above) + [magenta outline] [Os\textsuperscript{VI}O\textsubscript{2}(OH)\textsubscript{4}]\textsuperscript{2-}. (3): [yellow outline] [Os\textsuperscript{VIII}O\textsubscript{4}(OH)\textsubscript{2}]\textsuperscript{2-} + [green outline] 2[Fe\textsuperscript{II}(CN)\textsubscript{6}]\textsuperscript{4-} + 2H\textsubscript{2}O ⟶ [magenta outline] [Os\textsuperscript{VI}O\textsubscript{2}(OH)\textsubscript{4}]\textsuperscript{2-} + 2[Fe\textsuperscript{III}(CN)\textsubscript{6}]\textsuperscript{3-} + 2OH\textsuperscript{-}. (4): [magenta outline] 2[Os\textsuperscript{VI}O\textsubscript{2}(OH)\textsubscript{4}]\textsuperscript{2-} ⟶ [yellow outline] [Os\textsuperscript{VIII}O\textsubscript{4}(OH)\textsubscript{2}]\textsuperscript{2-} + [black fill] Os\textsuperscript{IV}O\textsubscript{2}(solid) + 2OH\textsuperscript{-}. As printed, reaction (4) is balanced in Os and charge but not in O and H (left 12 O, 8 H; right 10 O, 4 H): two H\textsubscript{2}O are missing on the right. Reactions (1)--(3) balance.}

\begin{center}
\includegraphics[width=0.95\linewidth,height=0.55\textheight,keepaspectratio]{figures/fig4.png}
\end{center}
\noindent \textbf{Figure 4:} Temperature- and duration-dependent outcomes of reduced-osmium staining. Left: OsO\textsubscript{4} (yellow squares) diffuse across cytoplasmic membranes and are gradually reduced over time (pink and black squares). Middle: ferrocyanide (green squares) reacts with ionized Os\textsuperscript{VIII} (yellow open squares) in the aqueous phase. Note that slow washout of water-soluble Os\textsuperscript{VI} species (grey arrows) in the presence of ferrocyanide. This may explain why the precipitate ring of rOTO is not found at the sample periphery, where the washout effect is particularly prominent. Right: short incubation (top left) yields uniform staining (black squares) for small-sized samples. In contrast, prolonged incubation required for large samples (top right) at room temperature (RT) may cause damage and extraction of stained tissues associated with overfixation. Right: this can be substantially alleviated by performing incubation at 4 °C or on ice (bottom), which decelerates osmium-related reactions and thereby preserves the integrity of cellular ultrastructure.
\alttext{Figure 4 is printed at the top of p.\,5, across the page. Schematic. A key in a black frame at upper left: a light-blue vertical bar with dark-blue edges = ``Membrane'' (blue text); a yellow filled square above a yellow open square = ``Os\textsuperscript{VIII}''; a pale pink filled square above a pale pink open square = ``Os\textsuperscript{VI}''; a green open hexagon = ``Fe\textsuperscript{II}''; a black filled square = ``Os\textsuperscript{IV}''. Left, ``Os\textsuperscript{VIII} incubation'': three vertical membranes between the vertical labels ``Sample surface'' (left) and ``Sample core'' (right); the membranes hold yellow filled squares (3, 2 and 2 from left to right) and one pink filled square each; yellow open squares sit in the gaps between membranes (one in each gap). Below, a dark-blue wedge, thick at the left and tapering to the right, labelled ``Passive diffusion''. An orange arrow points right to the middle panel, ``Fe\textsuperscript{II} incubation'': three membranes, each holding one black square and one pink filled square; green hexagons (3 in the left gap, 2 in the right gap) and pink open squares (2 in each gap) between the membranes. Below this panel two grey arrows bend down and then to the left (a thick one from under the left membrane, a thinner, lighter one from the right gap), labelled ``Washout'' in light grey. From the middle panel an orange arrow splits into two branches. Upper branch, ``Short incubation @RT'': three intact membranes with only black squares (2, 3 and 1). From there a further orange arrow ``Long incubation @RT'' leads to three membranes that are broken into shorter pieces, still with black squares (left: one short piece with 2; middle: two pieces with 1 each; right: one piece with 2). Lower branch, a long orange arrow ``Long incubation, @4°C'' to three intact full-length membranes with black squares (3, 3 and 2).}
\decompnote{The caption says ``ferrocyanide (green squares)''; in the figure and its key, Fe\textsuperscript{II} is drawn as green hexagons. ``Right:'' is printed twice in the caption.}

By and large, the proposed reaction model is consistent with experimental observations obtained across a range of osmium staining methods [26]. As for the original reduced osmium recipe, ferrocyanide rapidly depletes available OsO\textsubscript{4} to generate water-soluble Os\textsuperscript{IV} species. Although this process leads to more OsO\textsubscript{2} deposition onto cellular structures, the resultant low OsO\textsubscript{4} concentration restricts the tissue penetration depth of staining. More importantly, this mechanism provides a plausible explanation for why sequential incubation with OsO\textsubscript{4} and ferrocyanide allows OsO\textsubscript{4} to diffuse across greater distances before being converted into tissue-immobile OsO\textsubscript{2} deposits (Figure 4). Under certain circumstances, however, the water-soluble intermediate osmium products (Os\textsuperscript{VI}) may remain stabilized for an extended duration and may be washed out of the tissue prior to forming OsO\textsubscript{2}, the final reduction product. Until recently, the model seems unable to explain the progressive stain washout observed during prolonged ferrocyanide incubation (>3 h) [28], [31]. To resolve this discrepancy, it has been suggested that extended ferrocyanide incubation removes Os\textsuperscript{VI} species in the extracellular matrix through forming stable coordination compounds, which reduces background signals and thereby enhancing membrane contrast [31]. Whether this pathway represents the primary mechanism of ferrocyanide-enhanced osmium staining remains controversial, considering that OsO\textsubscript{4}-ferrocyanide mixtures do effectively stain various tissues within a short period and yield greater heavy-metal impregnation than OsO\textsubscript{4} treatment alone [20], [23]. Building on recent advances in low-temperature staining of insect [30], [47] and mammalian brain samples [31], we speculate that slow extraction of membrane lipids and biomolecules owing to osmium overfixation may largely account for osmium washout during prolonged incubation steps (Figure 4). This hypothesis is supported by cumulative evidence as well as our own experience showing that \pdfpage{6}prolonged OsO\textsubscript{4} incubation at room temperature induces severe tissue extraction and structural damage, whereas same incubation duration at 4 °C substantially alleviate such artifacts (Figure 5). Furthermore, osmium destaining predominantly occurs in peripheral sample regions, which is consistent with the time-dependent nature of overfixation, as outer tissue areas experience longer fixative exposure.
\decompnote{``water-soluble Os\textsuperscript{IV} species'' is printed so; elsewhere the paper calls the water-soluble intermediate Os\textsuperscript{VI} and gives Os\textsuperscript{IV} as the insoluble OsO\textsubscript{2} (Figure 3 and its caption).}

\begin{center}
\includegraphics[width=0.8\linewidth,height=0.55\textheight,keepaspectratio]{figures/fig5.png}
\end{center}
\noindent \textbf{Figure 5:} Evidence for temperature-dependent tissue extraction during staining. (A) brain samples were incubated in 2 \% buffered OsO\textsubscript{4} solutions at either room temperature (RT) or 4 °C. Low-temperature incubation markedly slows the darkening of the staining buffer, indicative of suppressed tissue extraction. Cacodylic buffer (CB) and phosphate buffer (PB) were used in this experiment. (B) SEM micrographs further verify that prolonged incubation at RT induces membranous structural damage (below), whereas low-temperature treatment effectively protects tissue ultrastructure (top).
\alttext{Figure 5 is printed at the top of p.\,6, across about two thirds of the page width. (A) Photographs, titled ``OsO\textsubscript{4} incubation''. Four rows of four capped plastic tubes, each with a small dark tissue piece at the bottom; rows labelled (vertical text, left) ``3 hours'', ``24 hours'', ``48 hours'', ``72 hours'', with an orange arrow pointing down along the left edge; columns labelled below ``CB @RT'', ``PB @RT'', ``CB @4°C'', ``PB @4°C''. 3 hours: the two RT tubes hold dark wine-red solution, the two 4°C tubes are almost colourless. 24 hours: RT tubes nearly black; 4°C tubes pale pink. 48 hours: RT tubes black; 4°C tubes medium rose. 72 hours: RT tubes black; 4°C tubes darker red, still translucent. Within each condition the CB and PB tubes look alike. (B) ``SEM image'': two greyscale SEM micrographs, each with a dark inset at upper left showing the whole sample block face (scale bar ``1 mm'') and a white scale bar ``1 μm'' at lower right. Top, labelled ``72 hours, @4°C'' (vertical): neuropil with continuous dark membranes, many small round and oval profiles, mitochondria with visible internal structure, larger pale profiles at left and upper right, and a thick black wavy band entering from the bottom centre; the inset block face is a fairly even mid-grey. Bottom, ``72 hours, @RT'': a uniform mottled grey texture with short dark streaks and specks in which membranes and profiles are hardly recognizable; only a few elongated spindle-shaped profiles with thick dark walls (upper left, right centre, left edge, top right) and a few small round profiles stand out; the inset block face is brighter in the middle with a darker rim.}

\section*{3 Future improvement of vEM sample preparation}

\noindent Although the preparation of large-volume tissue samples is technically feasible via prolonged low-temperature incubation, further protocol optimization with shortened processing durations and incorporate quality-control checkpoints in between critical steps is desirable to improve overall staining success rates. Since both stain-diffusion and chemical-reaction rates decline at low temperatures, incubations required for large specimens must be substantially prolonged. If overfixation-induced osmium washout cannot be fully suppressed even under low-temperature conditions, one would expect that excessively prolonged incubation may lead to noticeable destaining before thorough stain penetration is accomplished in oversized samples. Further optimization of chemical fixation is therefore required to minimize membrane lipid extraction, for example by microwave-assisted osmium fixation strategies [52].

Another unresolved challenge in vEM sample preparation lies in tissue fixation artifacts introduced by conventional chemical crosslinking of biomolecules. Previous studies have documented prominent morphological alternations in chemically fixed neural tissues, including reduced extracellular spaces [29], distorted neurite and synapse morphology [53], [54], and displaced synaptic vesicles [55]. The extent to which these fixation-induced morphological changes would influence biophysical modeling and connectome-based emulation remains to be clarified. As a promising approach, high-pressure-freezing (HPF) enables near-native ultrastructural preservation for tissue specimens up to 200 μm thick by suppressing ice crystal formation through rapid vitrification under high pressure. In principle, water molecules within HPF-preserved tissues can be gradually substituted with acetone under cryo-condition (freeze substitution), allowing subsequent resin embedding at room temperature. Afterwards, these samples can be handled normally for both serial section-based and serial block-face-based vEM (Figure 6), except for lower \pdfpage{7}contrast compared to rOTO-treated samples. Unfortunately, the poor solubility of ferrocyanide and TCH in anhydrous acetone makes existing \emph{en bloc} staining protocols difficult to implement during or after freeze substitution. To fully unlock the potential of HPF-preserved tissues for large-scale vEM applications, future methodological advances are needed to establish effective \emph{en bloc} staining strategies compatible with the freeze-substitution workflow. Promising approaches include the use of an acetone-based cocktail solution containing OsO\textsubscript{4}, uranyl acetate, tannic acid, and potassium permanganate [56]. Similarly, pyrogallol and gallic acid are expected to serve as potential osmiophilic candidates for stain amplification [27], [57], given their high solubility in acetone. Alternatively, supplementing freeze-substitution solutions with 5 \% water has been proposed to enable the use of water-soluble stains for cellular ultrastructure and membrane, such as a mixture of permanganate and uranium salts, as a cost-effective replacement for OsO\textsubscript{4} [58]. More recently, complete sample rehydration is carried out after freeze substitution using an acetone solution containing 2 \% OsO\textsubscript{4} and 3 \% water. This workflow, termed FSaqOTO [59], supports the subsequent aqueous osmium amplification procedure as well as post-staining with uranium and lead solutions. Furthermore, emerging image processing tools [60], [61], originally developed to boost vEM imaging throughput, can be repurposed to compensate for the reduced image contrast typically observed in HPF-preserved samples (Supplementary Video 1).

\begin{center}
\includegraphics[width=0.95\linewidth,height=0.55\textheight,keepaspectratio]{figures/fig6.png}
\end{center}
\noindent \textbf{Figure 6:} Block-face imaging of HPF-preserved brain sample. Left: example SEM micrograph of a freshly dissected mouse cortical sample preserved through HPF. The sample was subsequently subjected to freeze substitution with acetone liquid containing osmium and uranium, followed by resin embedding. The block-face was normally prepared for SEM imaging. Right: high-magnification view of the region outlined by the white box in the left panel. Axon viscosity (red solid arrow) and extracellular space (red open arrow) are clearly visualized. Beam current: 0.8 nA, voltage: 2.0 keV, dwell time: 1.5 μs.
\alttext{Figure 6 is printed at the top of p.\,7, across the page. Two greyscale SEM block-face images side by side, each with a black scale bar below its lower right corner. Left (wide, about 1.5:1; scale bar ``10 μm''): low-magnification view of mouse cortex; grey neuropil densely speckled with small black dots and small pale round profiles, a few larger pale cell bodies with darker nuclei (upper left, upper centre, upper right, centre), elongated pale vessel-like profiles in several directions, and a small dark elongated cell at lower left; a small white square box at upper left (about one fifth from the left edge) marks the region shown at right. Right (about 1.3:1; scale bar ``1 μm''): neuropil with thin dark continuous membranes outlining many round and oval profiles of different sizes, mitochondria, and several solid black round or oval profiles; pale, almost white gaps between some profiles. A red filled arrowhead (lower centre) points right at a dark, round profile lying inside a larger round profile; a red open arrowhead (right of centre) points down and to the left at a pale, almost white space between profiles, next to a small black dot.}

Last but not least, multiplexed labeling and correlative imaging with light microscopy (LM) represent important trends in modern vEM workflows [62]–[65]. Recent advances have enabled cell- and organelle-specific labeling using APEX-based probes [66], [67] as well as fixation-resistant and chemically stable fluorescent proteins [68]–[70], allowing for identification of genetically targeted cellular structures under EM. To date, however, these strategies have only been validated in samples processed with a single round of osmification [69] or with uranyl acetate to replace strong oxidative osmium compounds [64]. Further optimization of stain deposition density is therefore required to achieve unambiguous discrimination between labeled and unlabeled cyto-structures, while maximally preserving fluorescence after enhanced osmification and resin embedding for \emph{in situ} correlative LM and vEM imaging.

\section*{4 Summary}

\noindent Optimized osmium-based \emph{en bloc} staining has made large-volume and even whole-brain mapping technically achievable via commercial high-throughput vEM. With the growing advancement of high-speed imaging techniques and artificial intelligence-empowered acquisition optimization, saturated membrane staining is anticipated to no longer be the primary bottleneck for vEM experiments. Instead, increasing research attention will shift toward other exciting directions: developing accelerated or simplified workflows that can accomplish sample preparation within hours to days, enhancing contrast for HPF-preserved samples for large-scale high-throughput vEM, and improving the \pdfpage{8}compatibility of \emph{en bloc} staining with versatile genetic labeling strategies.

\bigskip
\noindent\textbf{Acknowledgments:} We thank Drs Mark Ellisman (UCSD) and Kun Song (PKU) for discussion and suggestion. We thank Bingyu Tan, Rui Ma, and Yinghong Lin (Shanghai NanoPort, Thermo Fisher Scientific) for support in image acquisition using FIB-SEM (Hydra Bio). We thank Bohao Chen (CAS-IA) for performing image denoising and Haoyu Wang (Shanghai Institute of Precision Medicine) for generating video illustration. This study was supported by the National Natural Science Foundation of China (32171461 to H.H.) and the Science and Technology Innovation 2030 (STI 2030) Major Projects (2021ZD0204500, 2021ZD0204503 to L.L.).\\
\textbf{Research ethics:} Not applicable.\\
\textbf{Informed consent:} Not applicable.\\
\textbf{Author contributions:} All authors have accepted responsibility for the entire content of this manuscript and approved its submission.\\
\textbf{Use of Large Language Models, AI and Machine Learning Tools:} None declared.\\
\textbf{Conflict of interest:} The authors state no conflict of interest.\\
\textbf{Research funding:} The National Natural Science Foundation of China (32171461 to H.H., 82171133 to Y.H.) and the Science and Technology Innovation 2030 (STI 2030) Major Projects (2021ZD0204500, 2021ZD0204503 to L.L.).\\
\textbf{Data availability:} Not applicable.

\section*{References}

\begin{hangpars}
\small
\pdfpage{8}[1] C. J. Peddie \emph{et al}., “Volume electron microscopy,” \emph{Nat. Rev. Methods Primers}, vol. 2, p. 51, 2022.

[2] A. A. Wanner, M. A. Kirschmann, and C. Genoud, “Challenges of microtome-based serial block-face scanning electron microscopy in neuroscience,” \emph{J. Microsc.}, vol. 259, no. 2, pp. 137−42, 2015.

[3] K. L. Briggman and D. D. Bock, “Volume electron microscopy for neuronal circuit reconstruction,” \emph{Curr. Opin. Neurobiol.}, vol. 22, no. 1, pp. 154−61, 2012.

[4] B. Titze and C. Genoud, “Volume scanning electron microscopy for imaging biological ultrastructure,” \emph{Biol. Cell}, vol. 108, no. 11, pp. 307−23, 2016.

[5] M. Helmstaedter, “Synaptic-resolution connectomics: Towards large brains and connectomic screening,” \emph{Nat. Rev. Neurosci.}, vol. 27, no. 2, pp. 101−20, 2026.

[6] R. Schalek \emph{et al}., “Development of high-throughput, high-resolution 3D reconstruction of large-volume biological tissue using automated tape collection ultramicrotomy and scanning electron microscopy,” \emph{Microsc. Microanal.}, vol. 17, no. S2, pp. 966−7, 2017.

[7] W. Yin \emph{et al}., “A petascale automated imaging pipeline for mapping neuronal circuits with high-throughput transmission electron microscopy,” \emph{Nat. Commun.}, vol. 11, no. 1, p. 4949, 2020.

[8] A. T. Kuan \emph{et al}., “Scalable electron tomography for connectomics,” \emph{bioRxiv}, 2024, \url{https://doi.org/10.1101/2024.06.05.597487}.

[9] T. Templier, “MagC, magnetic collection of ultrathin sections for volumetric correlative light and electron microscopy,” \emph{Elife}, vol. 8, 2019, \url{https://doi.org/10.7554/elife.45696}.\localcopy{templier2019}

[10] K. A. Fulton, P. V. Watkins, and K. L. Briggman, “GAUSS-EM, guided accumulation of ultrathin serial sections with a static magnetic field for volume electron microscopy,” \emph{Cell Rep. Methods}, vol. 4, no. 3, 2024, Art. no. 100720.

[11] N. Medina \emph{et al}., “DRIFT-EM enables direct wafer retrieval of ultrathin serial sections for large-volume electron microscopy,” \emph{Cell Rep. Methods}, vol. 6, no. 6, 2026, Art. no. 101429.

[12] W. Denk and H. Horstmann, “Serial block-face scanning electron microscopy to reconstruct three-dimensional tissue nanostructure,” \emph{PLoS Biol.}, vol. 2, no. 11, 2004, Art. no. e329.

[13] G. Knott, H. Marchman, D. Wall, and B. Lich, “Serial section scanning electron microscopy of adult brain tissue using focused ion beam milling,” \emph{J. Neurosci.}, vol. 28, no. 12, pp. 2959−64, 2008.

[14] K. J. Hayworth \emph{et al}., “Gas cluster ion beam SEM for imaging of large tissue samples with 10 nm isotropic resolution,” \emph{Nat. Methods}, vol. 17, no. 1, pp. 68−71, 2020.

[15] H. Ma \emph{et al}., “PIBE: Parallel ion beam etching of sections collected on wafer for ultra large-scale connectomics,” \emph{BioRxiv}, 2025, \url{https://doi.org/10.1101/2025.04.25.650569}.

[16] H. B. Nguyen \emph{et al}., “Conductive resins improve charging and resolution of acquired images in electron microscopic volume imaging,” \emph{Sci. Rep.}, vol. 6, no. 1, 2016, Art. no. 23721.\localcopy{nguyen2016}

[17] Y. Kubota \emph{et al}., “A carbon nanotube tape for serial-section electron microscopy of brain ultrastructure,” \emph{Nat. Commun.}, vol. 9, no. 1, p. 437, 2018.\localcopy{kubota2018}

[18] T. J. Deerinck, T. Shone, E. Bushong, R. Ramachandra, S. Peltier, and M. Ellisman, “High-performance serial block-face SEM of nonconductive biological samples enabled by focal gas injection-based charge compensation,” \emph{J. Microsc.}, vol. 270, no. 2, pp. 142−9, 2018.

[19] J. W. Lichtman and W. Denk, “The big and the small: Challenges of imaging the brain’s circuits,” \emph{Science}, vol. 334, no. 6056, pp. 618−23, 2011.

[20] D. L. White, J. E. Mazurkiewicz, and R. J. Barrnett, “A chemical mechanism for tissue staining by osmium tetroxide-ferrocyanide mixtures,” \emph{J. Histochem. Cytochem.}, vol. 27, no. 7, pp. 1084−91, 1979.\localcopy{white1979}

[21] A. M. Seligman, H. L. Wasserkrug, and J. S. Hanker, “A new staining method (OTO) for enhancing contrast of lipid − containing membranes and droplets in osmium tetroxide − fixed tissue with osmiophilic thiocarbohydrazide(TCH),” \emph{J. Cell Biol.}, vol. 30, no. 2, pp. 424−32, 1966.\localcopy{seligman1966}

[22] M. C. Willingham and A. V. Rutherford, “The use of osmium-thiocarbohydrazide-osmium (OTO) and ferrocyanide-reduced osmium methods to enhance membrane contrast and preservation in cultured cells,” \emph{J. Histochem. Cytochem.}, vol. 32, no. 4, pp. 455−60, 1984.\localcopy{willingham1984}

[23] J. C. Tapia \emph{et al}., “High-contrast en bloc staining of neuronal tissue for field emission scanning electron microscopy,” \emph{Nat. Protoc.}, vol. 7, no. 2, pp. 193−206, 2012.\localcopy{tapia2012}

[24] T. J. Deerinck, E. Bushong, V. Lev-Ram, X. Shu, R. Tsien, and M. Ellisman, “Enhancing serial block-face scanning electron microscopy to enable high resolution 3-D nanohistology of cells and tissues,” \emph{Microsc. Microanal.}, vol. 16, no. S2, pp. 1138−9, 2010.

\pdfpage{9}[25] K. L. Briggman, M. Helmstaedter, and W. Denk, “Wiring specificity in the direction-selectivity circuit of the retina,” \emph{Nature}, vol. 471, no. 7337, pp. 183−8, 2011.

[26] Y. Hua, P. Laserstein, and M. Helmstaedter, “Large-volume en-bloc staining for electron microscopy-based connectomics,” \emph{Nat. Commun.}, vol. 6, no. 1, p. 7923, 2015.\localcopy{hua2015}

[27] S. Mikula and W. Denk, “High-resolution whole-brain staining for electron microscopic circuit reconstruction,” \emph{Nat. Methods}, vol. 12, no. 6, pp. 541−6, 2015.\localcopy{mikula2015}

[28] S. Stroh, E. W. Hammerschmith, D. W. Tank, H. S. Seung, and A. A. Wanner, “In situ X-ray-assisted electron microscopy staining for large biological samples,” \emph{Elife}, vol. 11, 2022, \url{https://doi.org/10.7554/elife.72147}.\localcopy{stroh2022}

[29] M. Pallotto, P. V. Watkins, B. Fubara, J. H. Singer, and K. L. Briggman, “Extracellular space preservation aids the connectomic analysis of neural circuits,” \emph{Elife}, vol. 4, 2015, \url{https://doi.org/10.7554/elife.08206}.\localcopy{pallotto2015}

[30] A. A. Polilov, A. A. Makarova, S. Pang, C. Shan Xu, and H. Hess, “Protocol for preparation of heterogeneous biological samples for 3D electron microscopy: A case study for insects,” \emph{Sci. Rep.}, vol. 11, no. 1, p. 4717, 2021.

[31] K. Song, Z. Feng, and M. Helmstaedter, “High-contrast en bloc staining of mouse whole-brain and human brain samples for EM-based connectomics,” \emph{Nat. Methods}, vol. 20, no. 6, pp. 836−40, 2023.\localcopy{song2023}

[32] N. Karlupia \emph{et al}., “Immersion fixation and staining of multicubic millimeter volumes for electron microscopy-based connectomics of human brain biopsies,” \emph{Biol. Psychiatry}, vol. 94, no. 4, pp. 352−60, 2023.\localcopy{karlupia2023}

[33] Y. Hua \emph{et al}., “Connectomic analysis of thalamus-driven disinhibition in cortical layer 4,” \emph{Cell Rep.}, vol. 41, no. 2, 2022, Art. no. 111476.

[34] Y. Hua \emph{et al}., “Electron microscopic reconstruction of neural circuitry in the cochlea,” \emph{Cell Rep.}, vol. 34, no. 1, 2021, Art. no. 108551.

[35] C. Genoud, B. Titze, A. Graff-Meyer, and R. W. Friedrich, “Fast homogeneous en bloc staining of large tissue samples for volume electron microscopy,” \emph{Front. Neuroanat.}, vol. 12, p. 76, 2018.\localcopy{genoud2018}

[36] X. Lu, X. Han, Y. Meirovitch, E. Sjöstedt, R. L. Schalek, and J. W. Lichtman, “Preserving extracellular space for high-quality optical and ultrastructural studies of whole mammalian brains,” \emph{Cell Rep. Methods}, vol. 3, no. 7, 2023, \url{https://doi.org/10.1016/j.crmeth.2023.100520}.

[37] A. Gour \emph{et al}., “Postnatal connectomic development of inhibition in mouse barrel cortex,” \emph{Science}, vol. 371, no. 6528, 2021, \url{https://doi.org/10.1126/science.abb4534}.

[38] J. S. Phelps \emph{et al}., “Reconstruction of motor control circuits in adult Drosophila using automated transmission electron microscopy,” \emph{Cell}, vol. 184, no. 3, pp. 759−74 e718, 2021.

[39] C. Bosch \emph{et al}., “Functional and multiscale 3D structural investigation of brain tissue through correlative in vivo physiology, synchrotron microtomography and volume electron microscopy,” \emph{Nat. Commun.}, vol. 13, no. 1, p. 2923, 2022.

[40] N. Karagulyan \emph{et al}., “Gating of hair cell Ca(2+) channels governs the activity of cochlear neurons,” \emph{Sci. Adv.}, vol. 11, no. 25, 2025, Art. no. eadu7898.

[41] N. N. Guan \emph{et al}., “A specialized spinal circuit for command amplification and directionality during escape behavior,” \emph{Proc. Natl. Acad. Sci. U. S. A.}, vol. 118, no. 42, 2021, \url{https://doi.org/10.1073/pnas.2106785118}.

[42] T. M. Nguyen \emph{et al}., “Structured cerebellar connectivity supports resilient pattern separation,” \emph{Nature}, vol. 613, no. 7944, pp. 543−9, 2023.

[43] J. A. Bae \emph{et al}., “Functional connectomics spanning multiple areas of mouse visual cortex,” \emph{Nature}, vol. 640, no. 8058, pp. 435−47, 2025.

[44] N. L. Turner \emph{et al}., “Reconstruction of neocortex: Organelles, compartments, cells, circuits, and activity,” \emph{Cell}, vol. 185, no. 6, pp. 1082−100 e1024, 2022.

[45] C. M. Schneider-Mizell \emph{et al}., “Inhibitory specificity from a connectomic census of mouse visual cortex,” \emph{Nature}, vol. 640, no. 8058, pp. 448−58, 2025.

[46] S. Loomba \emph{et al}., “Connectomic comparison of mouse and human cortex,” \emph{Science}, vol. 377, no. 6602, 2022, Art. no. eabo0924.

[47] Z. Lu \emph{et al}., “En bloc preparation of Drosophila brains enables high-throughput FIB-SEM connectomics,” \emph{Front. Neural Circ.}, vol. 16, 2022, Art. no. 917251.

[48] C. Bosch \emph{et al}., “Nondestructive X-ray tomography of brain tissue ultrastructure,” \emph{Nat. Methods}, vol. 22, no. 12, pp. 2631−8, 2025.

[49] A. T. Kuan \emph{et al}., “Dense neuronal reconstruction through X-ray holographic nano-tomography,” \emph{Nat. Neurosci.}, vol. 23, no. 12, pp. 1637−43, 2020.

[50] G. Wildenberg \emph{et al}., “Photoemission electron microscopy for connectomics,” \emph{Proc. Natl. Acad. Sci. U. S. A.}, vol. 122, no. 48, 2025, Art. no. e2521349122.\localcopy{wildenberg2025}

[51] R. Li \emph{et al}., “OsO2 as the contrast-generating chemical species of osmium-stained biological tissues in electron microscopy,” \emph{ChemBioChem}, vol. 25, no. 20, Art. no. e202400311.
\decompnote{Reference [51] is printed without a year.}\localcopy{li2024}

[52] J. Hanker and B. Giammara, “Microwave-accelerated cytochemical stains for the image analysis and the electron microscopic examination of light microscopy diagnostic slides,” \emph{Scanning}, vol. 15, no. 2, pp. 67−80, 1993.

[53] H. Tamada, J. Blanc, N. Korogod, C. C. Petersen, and G. W. Knott, “Ultrastructural comparison of dendritic spine morphology preserved with cryo and chemical fixation,” \emph{Elife}, vol. 9, 2020, \url{https://doi.org/10.7554/elife.56384}.

[54] J. M. Griswold \emph{et al}., “Membrane mechanics dictate axonal pearls-on-a-string morphology and function,” \emph{Nat. Neurosci.}, vol. 28, no. 1, pp. 49−61, 2025.

[55] N. Korogod, C. C. Petersen, and G. W. Knott, “Ultrastructural analysis of adult mouse neocortex comparing aldehyde perfusion with cryo fixation,” \emph{Elife}, vol. 4, 2015, \url{https://doi.org/10.7554/elife.05793}.

[56] J. Guo \emph{et al}., “An optimized approach using cryofixation for high-resolution 3D analysis by FIB-SEM,” \emph{J. Struct. Biol.}, vol. 212, no. 1, 2020, Art. no. 107600.

[57] V. B. Wigglesworth, “A simple method for staining mitochondria,” \emph{Nature}, vol. 179, no. 4568, pp. 1033−4, 1957.

[58] M. Schauflinger, T. Bergner, G. Neusser, C. Kranz, and C. Read, “Potassium permanganate is an excellent alternative to osmium tetroxide in freeze-substitution,” \emph{Histochem. Cell Biol.}, vol. 157, no. 4, pp. 481−9, 2022.

[59] S. Belanger \emph{et al}., “A versatile enhanced freeze-substitution protocol for volume electron microscopy,” \emph{Front. Cell Dev. Biol.}, vol. 10, 2022, Art. no. 933376.

[60] Y. Meirovitch \emph{et al}., “SmartEM: Machine learning-guided electron microscopy,” \emph{Nat. Methods}, vol. 23, no. 1, pp. 193−204, 2026.

\pdfpage{10}[61] B. Chen \emph{et al}., “Volumetric denoising enables high-throughput volume electron microscopy and efficient downstream analysis,” \emph{Cell Rep. Methods}, 2026, Art. no. 101543, \url{https://doi.org/10.1016/j.crmeth.2026.101543}.

[62] L. M. Collinson, E. C. Carroll, and J. P. Hoogenboom, “Correlating 3D light to 3D electron microscopy for systems biology,” \emph{Curr. Opin. Biomed. Eng.}, vol. 3, pp. 49−55, 2017.

[63] N. Shafiei \emph{et al}., “Correlative light and electron microscopy for human brain and other biological models,” \emph{Nat. Protoc.}, vol. 20, no. 10, pp. 2994−3023, 2025.

[64] P. Ronchi \emph{et al}., “High-precision targeting workflow for volume electron microscopy,” \emph{J. Cell Biol.}, vol. 220, no. 9, 2021, \url{https://doi.org/10.1083/jcb.202104069}.

[65] D. P. Hoffman \emph{et al}., “Correlative three-dimensional super-resolution and block-face electron microscopy of whole vitreously frozen cells,” \emph{Science}, vol. 367, no. 7475, 2020, \url{https://doi.org/10.1126/science.aaz5357}.

[66] Q. Zhang, W. A. Lee, D. L. Paul, and D. D. Ginty, “Multiplexed peroxidase-based electron microscopy labeling enables simultaneous visualization of multiple cell types,” \emph{Nat. Neurosci.}, vol. 22, no. 5, pp. 828−39, 2019.

[67] S. S. Lam \emph{et al}., “Directed evolution of Apex2 for electron microscopy and proximity labeling,” \emph{Nat. Methods}, vol. 12, no. 1, pp. 51−4, 2015.

[68] Z. Fu \emph{et al}., “mEosEM withstands osmium staining and epon embedding for super-resolution CLEM,” \emph{Nat. Methods}, vol. 17, no. 1, pp. 55−8, 2020.

[69] H. Xiong \emph{et al}., “A highly stable monomeric red fluorescent protein for advanced microscopy,” \emph{Nat. Methods}, vol. 22, no. 6, pp. 1288−98, 2025.

[70] M. G. Paez-Segala \emph{et al}., “Fixation-resistant photoactivatable fluorescent proteins for CLEM,” \emph{Nat. Methods}, vol. 12, no. 3, pp. 215−8, 214 p following 218, 2015. \url{https://doi.org/10.1038/nmeth.3225}.
\end{hangpars}

\bigskip
\noindent \textbf{Supplementary Material:} This article contains supplementary material (https://doi.org/10.1515/mim-2026-0015).

\end{document}
